\section{Conclusion and Future Work}

We presented a knowledge-graph and rule co-optimization framework in which rule acquisition unlocks executable graph repairs and the value of further rule generation depends on the current graph state. The resulting implementation is a genuine Double-DQN system with replay, a target network, masked feasible actions, and archived state transitions.

On ten paired controlled-corruption runs, Double DQN reaches normalized final quality $0.9812\pm0.0013$ and AUC $0.9531\pm0.0021$. It improves final quality by 0.0065 and AUC by 0.0166 over the strongest non-oracle fixed schedule, with $p=0.000977$ in both paired one-sided Wilcoxon tests, while reducing rule-generation calls from eight to four. Standard DQN has a similar mean, so the evidence supports learned state-conditioned control; it does not establish a mean advantage of Double DQN over DQN. The model-informed myopic policy remains a slightly stronger upper bound.

The rule experiments show complementary behavior at two levels. In the full archived log, the dual strategy contributes 71.9\% more canonicalized candidates than the stronger single strategy. On the designed RuleTest-94 suite, deletion and augmentation families detect disjoint groups and their union covers all included cases. This suite result establishes executable coverage of known families rather than universal precision.

The next evaluation step is to replace controlled mutations with independently collected real repair trajectories and to create held-out domain test sets with per-case predictions from rule-mining baselines. Gold triple annotations are also needed for the API extraction benchmark, whose current category and keyword measures are silver indicators. These additions would test whether the learned scheduling benefit and generated constraints transfer beyond the controlled environment used here.
